% !TEX program = pdflatex
%
% Cornell-style meeting notes template.
%
% Layout (per page):
%   +----------------------------------------------------+
%   |            Header: meeting details & metadata      |
%   +----------------------------------------+-----------+
%   |                                        |   Cue      |
%   |           Main panel for notes         |  column    |
%   |                                        |  (hand-    |
%   |                                        |  writing)  |
%   +----------------------------------------+-----------+
%   |           Summary band (for hand writing)            |
%   +----------------------------------------------------+
%
% Usage: call \cornellpage{notes} once per page (see example at bottom).
% Header fields come from notes.yaml and apply to every page; page
% numbers are automatic. Adjustable knobs are in the
% "TEMPLATE CONFIGURATION" block below.
%
% Main content frame: written in notes.md and converted to
% build/cornell-content.tex (one \cornellFlow{...} per section) by
% `python3 scripts/markdown_to_pages.py` (or `make build`, which runs it
% automatically). \cornellFlow measures its content and automatically
% breaks it across as many pages as it needs. To force an extra break
% between two sections regardless, put a line containing only
%     <!-- pagebreak -->
% between them in notes.md.

\documentclass[11pt]{article}

\usepackage[T1]{fontenc}
\usepackage[letterpaper,margin=0.5in]{geometry}
\usepackage{tikz}
\usepackage{xcolor}
\usetikzlibrary{calc}
% Support for notes.md/pandoc output: task-list checkboxes (amssymb's
% \square/\boxtimes) and pipe tables (longtable/booktabs). All three are
% already pulled in by texlive-latex-recommended, so this costs nothing
% in image size.
\usepackage{amssymb}
\usepackage{longtable}
\usepackage{booktabs}
% Real clickable links/bookmarks for notes.md images and linked documents
% (see assets/); loaded last, as is customary for hyperref. hidelinks
% keeps the printed page free of colored/boxed link decoration.
\usepackage[hidelinks]{hyperref}

\pagestyle{empty}
\setlength{\parindent}{0pt}

% pandoc's LaTeX output uses \tightlist in itemize/enumerate bodies; it's
% normally defined by pandoc's own document template, which we bypass
% since we only take fragment output (see scripts/markdown_to_pages.py).
\providecommand{\tightlist}{%
  \setlength{\itemsep}{0pt}\setlength{\parskip}{0pt}}

% Bulleted list for cue-column entries (scripts/markdown_to_pages.py's
% "^<page> <text>" directive; see \cnCueText below). itemize's default
% leftmargin (~2.5em) eats too much of the narrow cue column, so this
% uses a plain `list` with a much smaller one instead, hugging the
% column's left edge.
\newenvironment{cnCueList}{%
  \begin{list}{\textbullet}{%
    \setlength{\leftmargin}{1em}\setlength{\itemindent}{0pt}%
    \setlength{\itemsep}{0pt}\setlength{\parsep}{0pt}%
    \setlength{\topsep}{0pt}\setlength{\partopsep}{0pt}%
  }%
}{%
  \end{list}%
}

% Numbered list for summary-band entries (scripts/markdown_to_pages.py's
% "^^<page> <text>" directive; see \cnSummaryText below). Same tight,
% narrow-column margin treatment as cnCueList above, but with an "N."
% label instead of a bullet -- summary-band text is typeset one column
% at a time (see \cnSummaryColW in \cornellpage), so it's just as narrow
% as the cue column.
\newcounter{cnSummaryListCtr}
\newenvironment{cnSummaryList}{%
  \begin{list}{\arabic{cnSummaryListCtr}.}{%
    \usecounter{cnSummaryListCtr}%
    \setlength{\leftmargin}{1.5em}\setlength{\itemindent}{0pt}%
    \setlength{\itemsep}{0pt}\setlength{\parsep}{0pt}%
    \setlength{\topsep}{0pt}\setlength{\partopsep}{0pt}%
  }%
}{%
  \end{list}%
}

% pandoc emits \st{...} for ~~strikethrough~~, normally provided by the
% `soul` package. Neither `soul` nor `ulem` is in the Docker image's
% minimal texlive set (both live only in texlive-latex-extra, a large
% pull), so this draws the strike line with tikz (already a dependency)
% instead. Unlike soul's \st, this treats its argument as one unbreakable
% box, so it won't line-wrap mid-phrase -- fine for a struck word or
% short phrase, not for striking out a whole paragraph.
\providecommand{\st}[1]{%
  \tikz[baseline=(cnStText.base)]{
    \node[inner sep=0pt,outer sep=0pt] (cnStText) {#1};
    \draw[line width=0.5pt]
      ($(cnStText.west)+(0,0.55ex)$) -- ($(cnStText.east)+(0,0.55ex)$);
  }%
}

% Images from notes.md (e.g. ![](assets/foo.png)) default-fit to the
% notes panel width instead of overflowing it.
\setkeys{Gin}{width=\linewidth}

% When notes.md sizes an image with only one dimension (e.g.
% ![](img.png){width=50pt}), pandoc pairs it with an explicit
% height=\textheight/\textwidth rather than leaving height unset. Without
% keepaspectratio, \includegraphics stretches to satisfy both dimensions
% exactly, badly distorting the image. Force keepaspectratio on every
% \includegraphics call so sized images scale instead of warping.
\let\cnOrigIncludegraphics\includegraphics
\renewcommand{\includegraphics}[2][]{\cnOrigIncludegraphics[#1,keepaspectratio]{#2}}

\newbox\cnFlowBox
\newbox\cnFlowSplit
\newbox\cnCueBox
\newbox\cnSummaryBox
\newbox\cnSummaryColA
\newbox\cnSummaryColB

% Directory the generated cornell-*.tex fragments below live in. The
% Makefile overrides this per invocation (via latexmk's -usepretex, which
% runs before this \providecommand, so it wins) with BUILDDIR -- e.g.
% build/example for `make build-example`, build/app for the Streamlit app
% -- so that concurrent/isolated builds actually use their own generated
% content instead of silently reading over each other's build/ directory.
\providecommand{\cnBuildDir}{build}

% ------------------------------------------------------------------
% Header fields sourced from notes.yaml.
% Run `python3 scripts/yaml_to_header.py` to (re)generate
% build/cornell-header.tex from notes.yaml before compiling; these
% \newcommand defaults keep the build working even if that hasn't been
% done yet.
% ------------------------------------------------------------------
\newcommand{\cnHdrTopic}{}
\newcommand{\cnHdrDate}{}
\newcommand{\cnHdrAttendees}{}
\newcommand{\cnHdrTime}{}
% Not typeset directly -- round-trips the app's IANA zone name (e.g.
% "America/Los_Angeles") through notes.yaml so its timezone dropdown can
% restore the exact zone on reload. The zone abbreviation shown on the
% page (e.g. "PDT") is baked into \cnHdrTime itself by the app.
\newcommand{\cnHdrTimezone}{}
% Typeset via \cnHdrTimeLine below, not appended to \cnHdrTime itself --
% keeps notes.yaml's "time" and "location" fields independent, so the
% Time row always reflects the current "location" value.
\newcommand{\cnHdrLocation}{}
\IfFileExists{\cnBuildDir/cornell-header.tex}{\input{\cnBuildDir/cornell-header.tex}}{}

% The Time row's displayed text: \cnHdrTime plus \cnHdrLocation, joined by
% ", " only when both are non-empty (mirrors the old single composed
% "time" string, but reconstructed from the two fields at build time
% instead of relying on the app to have concatenated them beforehand).
\newcommand{\cnEmptyMarker}{}
\newcommand{\cnHdrTimeLine}{%
  \ifx\cnHdrTime\cnEmptyMarker
    \cnHdrLocation%
  \else
    \ifx\cnHdrLocation\cnEmptyMarker
      \cnHdrTime%
    \else
      \cnHdrTime, \cnHdrLocation%
    \fi
  \fi
}

% Cue-column text for a given page number, from notes.md's "^<page> <text>"
% directive lines (see scripts/markdown_to_pages.py). \cnCueText{3} expands
% to whatever notes.md targeted at page 3, or nothing if none did.
\newcommand{\cnCueText}[1]{}
\IfFileExists{\cnBuildDir/cornell-cue.tex}{\input{\cnBuildDir/cornell-cue.tex}}{}

% Summary-band text for a given page number, from notes.md's
% "^^<page> <text>" directive lines (see scripts/markdown_to_pages.py).
\newcommand{\cnSummaryText}[1]{}
\IfFileExists{\cnBuildDir/cornell-summary.tex}{\input{\cnBuildDir/cornell-summary.tex}}{}

% ------------------------------------------------------------------
% TEMPLATE CONFIGURATION
% ------------------------------------------------------------------
\newlength{\cnPageW}       % drawable width  (slightly under \textwidth)
\newlength{\cnPageH}       % drawable height (slightly under \textheight)
\newlength{\cnHeaderHeight}
\newlength{\cnSummaryHeight} % the band at the bottom of the page
\newlength{\cnBodyTop}     % y-coordinate of the body/header boundary
\newlength{\cnCueWidth}
\newlength{\cnCueTextX}    % x-coordinate of the cue column's text area
\newlength{\cnCueTextW}    % usable text width of the cue column
\newlength{\cnSummaryColW} % usable text width of one summary-band column
\newlength{\cnSummaryColH} % usable text height of one summary-band column
\newlength{\cnSummaryColBX} % x-coordinate of the summary band's second column
\newlength{\cnNotesWidth}
\newlength{\cnHalfWidth}
\newlength{\cnPad}
\newlength{\cnBorderInset} % safety margin so thick-stroke ink stays inside \textwidth/\textheight
\newlength{\cnNotesPanelH} % usable height of the main notes panel, for \vsplit pagination

% Header/cue-column/summary-band sizing, as fractions of the drawable page.
\newcommand{\cnHeaderHeightFrac}{0.12}
\newcommand{\cnCueWidthFrac}{0.25}
\newcommand{\cnSummaryHeightFrac}{0.18}

\setlength{\cnPad}{0.12in}                   % inner padding for text
\setlength{\cnBorderInset}{2pt}

% ------------------------------------------------------------------
% Page geometry, header/cue-column/summary-band fractions, and spacing
% sourced from settings/page.yaml. Run `python3 scripts/yaml_to_settings.py`
% to (re)generate build/cornell-page-settings.tex from it before
% compiling; the defaults just above keep the build working even if
% that hasn't been done yet.
% ------------------------------------------------------------------
\IfFileExists{\cnBuildDir/cornell-page-settings.tex}{\input{\cnBuildDir/cornell-page-settings.tex}}{}

% \textwidth/\textheight are only finalized by geometry at
% \begin{document}, so the fractional lengths above must be computed
% then, not in the preamble (where they'd still hold article's defaults).
\AtBeginDocument{%
  \setlength{\cnPageW}{\textwidth}\addtolength{\cnPageW}{-\cnBorderInset}
  \setlength{\cnPageH}{\textheight}\addtolength{\cnPageH}{-\cnBorderInset}

  \setlength{\cnHeaderHeight}{\cnHeaderHeightFrac\cnPageH}    % header band
  \setlength{\cnSummaryHeight}{\cnSummaryHeightFrac\cnPageH}  % summary band
  \setlength{\cnBodyTop}{\cnPageH}
  \addtolength{\cnBodyTop}{-\cnHeaderHeight}

  \setlength{\cnCueWidth}{\cnCueWidthFrac\cnPageW}        % right-hand column
  \setlength{\cnNotesWidth}{\cnPageW}
  \addtolength{\cnNotesWidth}{-\cnCueWidth}

  \setlength{\cnCueTextX}{\cnNotesWidth}\addtolength{\cnCueTextX}{\cnPad}
  \setlength{\cnCueTextW}{\cnCueWidth}\addtolength{\cnCueTextW}{-2\cnPad}

  \setlength{\cnHalfWidth}{0.5\cnPageW}        % summary-band column width

  % Summary band splits into two side-by-side columns (left then right)
  % so a lot of "^^<page>" text grows sideways instead of vertically past
  % the band -- see \cornellpage below, which \vsplits the typeset text
  % at \cnSummaryColH into these two columns. Both columns are inset from
  % the page edges and the shared \cnHalfWidth divider by \cnPad.
  \setlength{\cnSummaryColW}{\cnHalfWidth}\addtolength{\cnSummaryColW}{-1.5\cnPad}
  \setlength{\cnSummaryColH}{\cnSummaryHeight}\addtolength{\cnSummaryColH}{-2\cnPad}
  \setlength{\cnSummaryColBX}{\cnHalfWidth}\addtolength{\cnSummaryColBX}{0.5\cnPad}

  % usable size of the main notes panel, for \vsplit-based pagination
  \setlength{\cnNotesTextW}{\cnNotesWidth}\addtolength{\cnNotesTextW}{-2\cnPad}
  \setlength{\cnNotesPanelH}{\cnBodyTop}
  \addtolength{\cnNotesPanelH}{-\cnSummaryHeight}
  \addtolength{\cnNotesPanelH}{-2\cnPad}
}

% scratch registers used while laying out each page
\newlength{\cnRowAY}
\newlength{\cnRowBY}
\newlength{\cnMeetW}
\newlength{\cnDateX}
\newlength{\cnDateW}
\newlength{\cnAttW}
\newlength{\cnTLX}
\newlength{\cnTLW}
\newlength{\cnCapMidX}
\newlength{\cnCapRightX}
\newlength{\cnNotesTextW}
\newlength{\cnNotesTextY}
\newlength{\cnSummaryTextY}    % y-coordinate of the summary band's text area

% A labelled fill-in blank: label at top-left (x,y); an optional typed
% value on the line below it, wrapped to width w so it can never spill
% into a neighboring field; and a rule underneath for writing by hand.
% x, y, w must be bare dimensions (no arithmetic).
\newcommand{\cnField}[5]{% x, y (top-left), width, label, value
  \cnFieldNoRule{#1}{#2}{#3}{#4}{#5}%
  \draw[gray!70] ($(#1,#2)+(0pt,-34pt)$) -- ($(#1,#2)+(#3,-34pt)$);
}

% Same as \cnField but without the rule underneath.
\newcommand{\cnFieldNoRule}[5]{% x, y (top-left), width, label, value
  \node[anchor=north west,font=\bfseries\footnotesize,inner sep=0pt]
    at (#1,#2) {#4};
  \node[anchor=north west,align=left,font=\footnotesize,inner sep=0pt,
        text width=#3] at ($(#1,#2)+(0pt,-11pt)$) {#5};
}

% ------------------------------------------------------------------
% \cornellpage{Notes content}
%
% Call once per page. Header fields (Topic/Date/Attendees/Time)
% come from \cnHdrTopic etc. (see notes.yaml) and are applied to every
% page automatically. Pass an empty group ({}) for "Notes content" to
% leave the main panel blank for handwriting instead of typed text.
% The page number in the header's upper-right corner is automatic.
% ------------------------------------------------------------------
\newif\ifcnfirstpage
\cnfirstpagetrue

\newcommand{\cornellpage}[1]{%
  % \global is required here: \cornellpage is always called from inside
  % a \cornellFlow (see below), and \newenvironment wraps every
  % environment's body in its own TeX group -- a plain (non-\global)
  % \cnfirstpagefalse would only last until that \cornellFlow's
  % \end{cornellFlow} closes the group, then silently revert to true.
  % With a single \cornellFlow block (one notes.md section, the common
  % case) that's invisible, since nothing rechecks it afterward; with a
  % second \cornellFlow block (from a "<!-- pagebreak -->" in notes.md),
  % its first \cornellpage call would wrongly see \cnfirstpage as true
  % again, skip \newpage, and every page from there on would render
  % with the previous page's already-typeset \thepage value.
  \ifcnfirstpage\global\cnfirstpagefalse\else\newpage\fi
  % --- precompute header field geometry -------------------------------
  \setlength{\cnRowAY}{\cnPageH}\addtolength{\cnRowAY}{-8pt}
  \setlength{\cnRowBY}{\cnRowAY}\addtolength{\cnRowBY}{-0.5\cnHeaderHeight}

  \setlength{\cnMeetW}{0.62\cnPageW}\addtolength{\cnMeetW}{-2\cnPad}
  \setlength{\cnDateX}{0.64\cnPageW}
  \setlength{\cnDateW}{0.36\cnPageW}\addtolength{\cnDateW}{-\cnPad}

  % Attendees/Time mirror Topic/Date's column split so Time lands
  % directly below Date.
  \setlength{\cnAttW}{\cnMeetW}
  \setlength{\cnTLX}{\cnDateX}
  \setlength{\cnTLW}{\cnDateW}

  \setlength{\cnCapMidX}{\cnHalfWidth}\addtolength{\cnCapMidX}{\cnPad}
  \setlength{\cnCapRightX}{\cnPageW}\addtolength{\cnCapRightX}{-\cnPad}

  \setlength{\cnNotesTextW}{\cnNotesWidth}\addtolength{\cnNotesTextW}{-2\cnPad}
  \setlength{\cnNotesTextY}{\cnBodyTop}\addtolength{\cnNotesTextY}{-\cnPad}
  \setlength{\cnSummaryTextY}{\cnSummaryHeight}\addtolength{\cnSummaryTextY}{-\cnPad}

  % Typeset this page's cue text.
  \setbox\cnCueBox=\vbox{\hsize=\the\cnCueTextW\linewidth=\hsize\footnotesize\noindent\cnCueText{\thepage}}%

  % Summary text is typeset at one column's width, then \vsplit at one
  % column's height: whatever fits lands in column 1, and any overflow
  % (rather than running past the band) lands in column 2 instead. A
  % page with little/no summary text ends up with an empty (\ifvoid)
  % column 2.
  \setbox\cnSummaryBox=\vbox{\hsize=\the\cnSummaryColW\linewidth=\hsize\footnotesize
    \noindent\cnSummaryText{\thepage}\par\vfil\penalty-10000}%
  \setbox\cnSummaryColA=\vsplit\cnSummaryBox to \cnSummaryColH\relax
  \ifvoid\cnSummaryBox
    % Nothing overflowed into column 2 -- clear it explicitly, since
    % otherwise it would still hold column 2's box from whichever
    % earlier page last had overflow (box registers are global and
    % persist across \cornellpage calls until reassigned).
    \setbox\cnSummaryColB=\hbox{}%
  \else
    \setbox\cnSummaryColB=\vsplit\cnSummaryBox to \cnSummaryColH\relax
  \fi

  \noindent
  \begin{tikzpicture}
    % --- outer border + row/column dividers (table skeleton) ---
    \draw[black,thick] (0,0) rectangle (\the\cnPageW,\the\cnPageH);
    \draw[black,thick] (0,\the\cnSummaryHeight) -- (\the\cnPageW,\the\cnSummaryHeight); % summary / body
    \draw[black,thick] (0,\the\cnBodyTop) -- (\the\cnPageW,\the\cnBodyTop);             % body / header
    \draw[black,thick] (\the\cnNotesWidth,\the\cnSummaryHeight) -- (\the\cnNotesWidth,\the\cnBodyTop); % notes / cue
    % No visible divider between the summary band's two columns (see
    % \cnHalfWidth/\cnSummaryColBX below) -- they're a layout/wrap detail,
    % not a boundary meant to be seen.

    % --- header fields (label + optional typed value) --------------------
    \cnField{\the\cnPad}{\the\cnRowAY}{\the\cnMeetW}{Topic:}{\cnHdrTopic}
    \cnField{\the\cnDateX}{\the\cnRowAY}{\the\cnDateW}{Date:}{\cnHdrDate}

    \cnFieldNoRule{\the\cnPad}{\the\cnRowBY}{\the\cnAttW}{Attendees:}{\cnHdrAttendees}
    \cnFieldNoRule{\the\cnTLX}{\the\cnRowBY}{\the\cnTLW}{Time:}{\cnHdrTimeLine}

    % --- region captions --------------------------------------------------
    \node[anchor=south west,font=\itshape\scriptsize,gray,inner sep=2pt]
      at (\the\cnPad,0pt) {};
    \node[anchor=south west,font=\itshape\scriptsize,gray,inner sep=2pt]
      at (\the\cnCapMidX,0pt) {};
    \node[anchor=south east,font=\itshape\scriptsize,gray,inner sep=2pt]
      at (\the\cnCapRightX,\the\cnSummaryHeight) {};

    % --- auto page number: upper-right corner of the header box ---------
    \node[anchor=north east,font=\itshape\scriptsize,gray,inner sep=2pt]
      at (\the\cnCapRightX,\the\cnPageH) {Page \thepage};

    % --- main notes panel (typed content, or blank for handwriting) -----
    \node[anchor=north west,align=left,
          text width=\the\cnNotesTextW,
          inner sep=0pt] at (\the\cnPad,\the\cnNotesTextY) {#1};

    % --- cue column text for this page (notes.md's "^<page> <text>") ----
    \node[anchor=north west,inner sep=0pt]
      at (\the\cnCueTextX,\the\cnNotesTextY) {\box\cnCueBox};

    % --- summary band text for this page (notes.md's "^^<page> <text>"),
    % column 1 then column 2 (overflow from column 1) ---
    \node[anchor=north west,inner sep=0pt]
      at (\the\cnPad,\the\cnSummaryTextY) {\box\cnSummaryColA};
    \node[anchor=north west,inner sep=0pt]
      at (\the\cnSummaryColBX,\the\cnSummaryTextY) {\box\cnSummaryColB};
  \end{tikzpicture}%
  \par
}

% ------------------------------------------------------------------
% \begin{cornellFlow} content \end{cornellFlow}
%
% Typesets content into the notes panel and, unlike \cornellpage, breaks
% it across as many pages as it actually needs: content is measured into
% a single tall box, then repeatedly \vsplit at the panel height, calling
% \cornellpage once per resulting chunk (each with the header and page
% number applied automatically, as usual).
%
% This is an environment rather than a \newcommand argument so that
% verbatim-family environments (pandoc's rendering of Markdown code
% blocks) can appear inside it: verbatim needs to read its body live,
% off the main input stream, which a macro argument's pre-scan would
% prevent.
%
% \cnfinalflow marks the very last cornellFlow block in the document
% (set by scripts/markdown_to_pages.py, which is the only thing that
% knows how many cornellFlow blocks notes.md produces); \cnlastpage is
% then true only for that block's own last \vsplit'd page -- i.e. the
% document's actual final page. \cnCueText/\cnSummaryText (see
% scripts/markdown_to_pages.py's render_directive_tex) check \cnlastpage
% to render a "^<page>"/"^^<page>" directive whose target page is beyond
% the document's real length there, instead of it silently never
% appearing.
% ------------------------------------------------------------------
\newif\ifcnfinalflow
\cnfinalflowfalse
\newif\ifcnlastpage
\cnlastpagefalse

\newenvironment{cornellFlow}{%
  \setbox\cnFlowBox=\vbox\bgroup
    \hsize=\cnNotesTextW
    \linewidth=\hsize
    \footnotesize
    \noindent
}{%
    \par\vfil\penalty-10000
  \egroup
  \cnFlowLoop
}
\def\cnFlowLoop{%
  \ifvoid\cnFlowBox
  \else
    \setbox\cnFlowSplit=\vsplit\cnFlowBox to \cnNotesPanelH\relax
    \cnlastpagefalse
    \ifvoid\cnFlowBox\ifcnfinalflow\cnlastpagetrue\fi\fi
    \cornellpage{\box\cnFlowSplit}%
    \cnFlowLoop
  \fi
}

\begin{document}

% ------------------------------------------------------------------
% One \cornellpage per page of notes.md (split on "<!-- pagebreak -->");
% see build/cornell-content.tex generation notes above. Each page
% automatically gets the same header (from notes.yaml) and the next
% page number.
% ------------------------------------------------------------------

\input{\cnBuildDir/cornell-content.tex}

\end{document}
